\clearpage
\section{Validation}

The validation strategy combines \textbf{unit tests} and an
\textbf{end-to-end RMI integration test} so that both local correctness and
distributed behavior are covered.

\begin{itemize}
    \item \texttt{GameStatusTest} validates the helper predicates used by the
    game-state model.
    \item \texttt{GameImplTest} validates move legality, turn alternation, win
    and draw detection, abandonment handling, snapshot consistency, and
    concurrent operations on a single match.
    \item \texttt{LobbyImplTest} validates room creation, room join, duplicate
    names, waiting-room listing, pruning of terminal games, and concurrent
    accesses to the lobby.
    \item \texttt{DtttRmiIntegrationTest} validates the complete distributed
    path through registry lookup, server binding, room creation, room join, and
    callback-based state propagation.
\end{itemize}

This coverage targets the main risks of the architecture. The unit tests verify
that the server remains the \textbf{single point of truth} for game rules and
state transitions, while the integration test verifies that remote discovery,
remote invocation, and client callbacks work together correctly in the actual
RMI-based execution model.

The validation is therefore not limited to rule correctness. It also checks
that requests are routed to the correct remote game instance and that
callback-driven updates expose consistent \texttt{BoardState} snapshots to
clients. Match isolation follows from the design choice of one
\texttt{GameImpl} instance per room, rather than from a dedicated multi-match
integration scenario in the current test suite.

The current validation scope is appropriate for the assignment: it demonstrates
functional correctness, distributed interaction correctness, and concurrency
control within the designed architecture. It does not aim to cover
\textbf{persistence}, \textbf{replicated servers}, or large-scale
\textbf{network-failure scenarios}, since these are outside the intended scope
of the project.
